\paragraph{Acyclic source dependence.}
A \emph{source order} is an order of the inputs in which the complete
predicate $\pi_i$ at source $i$ depends only on predecessors of $i$.
Equivalently, the graph recording dependence of one source predicate on
another input's passive status is acyclic. To see the equivalence, erase
later atomic statuses one at a time; every intermediate coloring remains
proper. Write $q|_{<i}$ for the predecessor restriction of a coloring.
The cofactor law below always concerns the entire function on the unfixed
inputs, including changes of an unfixed menu to a different atomic color.

\begin{lemma}[Captured-input rectangles]\label{core:lem:captured-rectangle}
Let $C$ be core on $W$, and let a source-reconstructing operation $F$
satisfy $C\subseteq F(X,v)\subseteq C\cup X$, with antitone source
predicate $P$ and the colored-cofactor independence law.
Let $R$ be a proper coloring, let $B\subseteq W$ be disjoint from its
colored inputs, and suppose $C$ captures every input in $B$ under $R$.
If $P(R)=1$, then the whole owner box
\[
 Q_{R,B}(E)=\bigwedge_{r\in I}[E_r\setminus B\subseteq R_r]
\]
is contained in $P$.
\end{lemma}
\begin{proof}
The passive source families of $F$ equal those of $C$. Consequently $F$
also captures every input in $B$ under $R$.
Take $E$ in the box and put $T=R\cup(E|_B)$. These colorings are
compatible, and $E\subseteq T$. If $T$ is proper, add its captured
inputs one at a time. At a probe for a missing owner, cofactor independence
preserves the sole $X$-winner, so $P(T)=P(R)=1$. Antitonicity gives
$P(E)=1$.

Suppose instead every owner occurs in $T$; some variables may still be
uncolored. Choose $\ell$ unused by $R$. Its entire block $T_\ell$
lies in $B$. The preceding proper-coloring argument gives
$P(T^{\widehat\ell})=1$.
For any owner $k$, suppose $P(T^{\widehat k})=0$.
The new-source cofactor law would make $F$ independent of $X$ after
fixing $T^{\widehat k}$. Extend this fixing by all of $T$, and give
every still-uncolored passive input $I\setminus\{\ell\}$.
All possible $\ell$-sources in $C$ are in $T_\ell$, and each is
captured already under $R$. Thus $C$ omits $\ell$.
With $X=I\setminus\{\ell\}$, $F$ omits it too; with $X=\{\ell\}$,
$P(T^{\widehat\ell})=1$ supplies it. This contradicts independence.
Hence $P(T^{\widehat k})=1$ for every $k$.
Choose $k$ unused by $E$. Then $E\subseteq T^{\widehat k}$, giving
$P(E)=1$ by antitonicity.
\end{proof}

\begin{lemma}[Reduction of a source pattern]\label{core:lem:source-pattern-reduction}
Let $C$ be core and have a source order.
For a proper coloring $q$, put
\[
 J(q)=\{i:\pi_i^C(q|_{<i})=1\},\qquad R(q)_r=q_r\cap J(q).
\]
Every input outside $J(q)$ is captured by $C$ after fixing $R(q)$.
\end{lemma}
\begin{proof}
Process inputs in source order. Previously excluded inputs have already
been captured under the retained predecessor colors. In a canonical
source-$i$ probe, choose an observer unused by $q$. Change the excluded
predecessors from neutral menus to their $q$-values. Cofactor independence
preserves the queried winner, so
\[
 \pi_i^C(R(q)|_{<i})=\pi_i^C(q|_{<i}).
\]
If this value is zero, input $i$ is captured under the retained prefix,
and remains captured after fixing the rest of $R(q)$. This proves the
induction.
\end{proof}

\begin{theorem}[Acyclic source compatibility and ordered spines]
\label{core:thm:acyclic-source-guards}
Fix finite $K\geq3$ and $n\geq1$ inputs. Let $F$ be a nonempty
operation reconstructed from its complete source predicates. Suppose
these predicates agree across observer overlaps, are antitone, and
admit a source order. Then $F$ is core if and only if a zero source
predicate makes its proper atomic cofactor independent of that source.

When these conditions hold, discard identically unused sources. For each
remaining source $i$ and each maximal positive proper predecessor coloring
$q$ of $\pi_i$, compute $J_i(q)=\{j<i:\pi_j(q|_{<j})=1\}$.
Start a spine at input $i$. Process $j\in J_i(q)$ in decreasing source
order, attaching input $j$ once by $\star_r$ for each owner $r$ with
$j\notin q_r$. Fix any order within each such group.
The union of these spines is exactly $F$.
With union named it has at most
$(1+K(n-1))nK^n$ leaves.
Moreover every singleton profile is principal:
\[
 \Phi_r(F)(A)=1\quad\Longleftrightarrow\quad M_r\subseteq A,
 \qquad M_r=\{i:\pi_i(\text{all predecessors colored }r)=1\}.
\]
\end{theorem}
\begin{proof}
Necessity is \cref{core:cor:colored-cofactor-dependence}.
We first prove an interval construction over an already core baseline $C$
with a source order. Let $F_P$ be a source-reconstructing interval candidate
above $C$ satisfying the same cofactor law. Inputs unused by $C$ are
also globally inessential in $F_P$ and may be omitted. The first remaining
input of $C$ is retained, so its nonempty prefixes are core neutral
residuals. Fix $P(q)=1$.
Set $J=J(q)$, $R=R(q)$, and $B=W\setminus J$.
Antitonicity gives $P(R)=1$. By the two preceding lemmas,
\[
 Q_q(E)=\bigwedge_{r\in I}[E_r\cap J\subseteq q_r]
 \quad\Longrightarrow\quad P(E)=1.
\]
The indicated spine at $X$, using the variables in $J$ in decreasing
order, has exactly $Q_q$ as its $X$-source predicate. At a probe the
passive inputs omit the queried winner, so only $X$ can supply it.
The forbidden-owner gates apply precisely the displayed factors; the
queried owner's factor is automatically true.

We show that this spine lies below $C\cup X$. A possible winner $k$
outside that bound must come from a particular helper occurrence of an
input $i\in J$. Thus $k\in v_i$ and $k\notin C(v)$.
Let $T$ be the foreign-$k$ atomic coloring of the predecessors of $i$.
Source decomposition and the source order give $\pi_i^C(T)=0$.
The first retained input cannot be such a helper, since $C$ contains it.
For any other $i$, consider the actual core prefix $C_{\leq i}$ as an
interval extension of $C_{<i}$ by source $i$. These prefixes are neutral
residuals of $C$; their earlier source predicates are unchanged.
Its source-$i$ predicate is positive at $q|_{<i}$.
Apply the captured-input rectangle lemma to this prefix extension.
Its earlier active set is exactly $J(q)\cap\{j:j<i\}$, since earlier
predicates are unchanged.
If $T$ obeyed the rectangle determined by $q|_{<i}$, it would force
$\pi_i^C(T)=1$, a contradiction. Therefore some earlier active
$j\in J$, $j<i$, has $v_j=\{r\}$ for $r\ne k$ and $j\notin q_r$.
The $j$-group is attached later on the spine. Its $r$-owned gate
absorbs to $\{r\}$ and kills this particular helper occurrence's
contribution of $k$. Every helper occurrence is covered by this argument;
any later reintroduction must come from another occurrence, covered in turn.
Thus no helper can supply a winner outside $C\cup X$.

Union with $C$ now realizes $Q_q$ in the interval. Union over maximal
positive $q$ realizes $P$: every positive coloring extends to a maximal
one, and each rectangle is antitone, contains its defining $q$, and lies
in $P$. The zero predicate uses $C$ alone.

For the global result, remove sources with identically zero predicates.
The cofactor law at the empty coloring makes them globally inessential.
The first remaining source has no predecessors and has predicate one,
so is retained and its prefix is a projection.
Every topological prefix is the winnerwise neutral residual deleting
later inputs: later inputs contribute no queried winner, while earlier
predicates ignore their statuses. The cofactor law inherits to these
prefixes by neutralizing the same later inputs in both compared valuations.
Consequently each prefix is an interval extension of the preceding one.
Induction with the interval construction makes every prefix core.
Expanding its union of spines gives exactly the stated normal form;
only original input variables occur in the helpers.

Each proper predecessor coloring is represented by a canonical probe
for one of its missing observers, with later inputs neutral. Thus the
number of chosen spines is at most the $nK^n$ source observations.
Each spine uses at most $1+K(n-1)$ leaves, proving the bound.

Finally, apply pattern reduction to the all-$r$ coloring.
Fixing $M_r$ to $r$ captures every other input, so forces the singleton
$r$. Conversely, if an $r$-forcing coalition $A$ omitted some $i\in M_r$,
antitonicity would give
$\pi_i(A|_{<i})\geq\pi_i(\text{all predecessors }r)=1$.
A sole-source probe at $i$ for any foreign winner would contradict the
forced singleton. The core singleton-profile factorization identifies
these forcing coalitions with $\Phi_r(F)$, giving the final formula.
\end{proof}
